\section{Actual PL weights on the same original normal proper ratio model}

\begin{definition}[The true whole normalized original Gauss lift]
\label{def:actual-original-normalized-simultaneous-gauss-point}
\lean{EnvelopeProperty.ProjSubvariety.normalizedSimultaneousGaussPoint}
\lean{EnvelopeProperty.ProjSubvariety.map_normalizedSimultaneousGaussPoint}
\leanok
\uses{def:actual-original-normalized-simultaneous-model,thm:actual-original-pl-variety-product-normal,thm:actual-normal-affine-neighborhood-over-open-isomorphism,def:actual-normalization-relative-point-lift}
On an original projectively normal variety, the actual simultaneous Gauss
point has a true whole lift to the same actual normalized model. Projection
retains the entire original Gauss point and its original residue valuation.
\end{definition}
\begin{proof}
Finiteness of the two original ideal logarithms puts the original Gauss
point off the actual simultaneous blowup center. The true whole original
blowup lift lies over that same open. The original product's actual stalks
are normal by projective normality, and the original blowup is an actual
isomorphism over that open. The proved true open-isomorphism stalk comparison
therefore constructs a normal affine neighborhood of the original blowup
support. The actual normalization is a whole isomorphism over that original
normal affine neighborhood; the proved actual relative-point lift through
that isomorphism gives the stated whole original point and projection identity.
\end{proof}

\begin{definition}[The actual normalized ratio line and generic section]
\label{def:actual-original-normalized-simultaneous-ratio-line}
\lean{EnvelopeProperty.ProjSubvariety.normalizedSimultaneousRatioLine}
\lean{EnvelopeProperty.ProjSubvariety.normalizedSimultaneousRatioLine_isInvertible}
\lean{EnvelopeProperty.ProjSubvariety.normalizedSimultaneousRatioSection}
\leanok
\uses{def:actual-original-normalized-simultaneous-gauss-point,def:actual-original-normalized-simultaneous-model,cor:actual-cartier-ratio-model-weight}
Pull the two same original Cartier ideals on the simultaneous blowup to its
actual normalization. Their actual ratio is an invertible original sheaf.
At the true normalized Gauss point take its actual canonical generic ratio
section: the first ideal's rational unit section tensored with the second
ideal inclusion's actual dual-transpose fiber section.
\end{definition}
\begin{proof}
The actual simultaneous blowup makes both original ideals Cartier. Its
normalization is integral and the projection is dominant, so the proved
dominant Cartier pullback theorem retains these actual Cartier ideals.
Actual tensor and dual evaluation give invertibility of their actual ratio.
Original ideal-norm pullback retains both entire ideal functions at the
constructed entire point. Thus the first norm is nonzero, as its original
Gauss logarithm is finite. The canonical original generic ratio construction
therefore supplies the stated actual section of this same actual sheaf.
\end{proof}

\begin{theorem}[True normalized ratio weight and exact original PL potential]
\label{thm:actual-original-normalized-simultaneous-ratio-weight}
\lean{EnvelopeProperty.ProjSubvariety.normalizedSimultaneousRatioWeight}
\lean{EnvelopeProperty.ProjSubvariety.normalizedSimultaneousRatioWeight_eq}
\lean{EnvelopeProperty.ProjSubvariety.IsPL.exists_normalizedSimultaneousRatioMetric_models}
\leanok
\uses{def:actual-original-normalized-simultaneous-ratio-line,thm:actual-original-normalized-simultaneous-model-proper,cor:actual-cartier-ratio-model-weight}
Over any original field, the negative logarithm of the actual model
norm of this true generic ratio section on the same original normal proper
model is exactly the difference of the two original intrinsic ideal
logarithms. Consequently any original PL function on a projectively normal
original variety admits a positive integer scale and integer constant for
which this actual original normalized ratio weight gives its exact value.
\end{theorem}
\begin{proof}
The actual normalized model is proper over its same polynomial base by
finite normalization over the original field and original field finite type. The true
canonical Cartier ratio norm formula therefore identifies the actual negative
logarithmic section norm with the two pulled ideal logarithms' difference.
True geometric ideal pullback through normalization, the entire normalized
point projection identity, and true geometric ideal pullback through the
original simultaneous blowup identify this with the original Gauss ideal-log
difference. Finally use the already proved original PL ideal-model producer:
it supplies the two actual ideals, their finite original Gauss logarithms,
positive integer scale, integer constant, and exact original ideal-difference
formula. Substitute the proved actual normal proper ratio-weight identity.
No normal model, metric compatibility, Gauss lift or numerical positivity
certificate is assumed. Relative nefness and the eventual original polarized
line comparison are further geometric conclusions, not asserted here.
\end{proof}
